\clearpage
\section{Concrete copying specimens and exact class semantics}
\label{sec:copying-specimens}
The combined copying UAM now supplies a concrete instance of the abstract
specimen interface. Its synchronized context-step law is proved from actual
recursive graph fillings. Its machine is deterministic on classes under
recursive graph isomorphism and depth padding, and every class execution
reflects to an execution from any chosen original representative with exactly
the same length. All existing binding-free program and term observations are
therefore preserved. The terms of Proposition~12.2, Figure~16's translation,
the complete operation alphabet and all rule premises are unchanged.

The locality argument follows the three conclusions of PDF Lemma~F.5
(HTML Lemma~39, \href{https://arxiv.org/html/1907.01257v5\#A6}{Appendix~F}):
the two filled states agree on finality; output closure prevents a jointly
rooted token from exiting a hole; and a settled specimen can synchronize a
step on both sides to a remote specimen. Here tokens are typed coordinates
on a shared focus-free context. Filling certificates recurse through original
box bodies and retain their complete ordered interfaces.

The proved dependency chain is: ordered incidence preservation and boundary
reflection; one-way operation-path obstruction; recursive paired fillings;
exact representation bisimulation and literal class determinism; synchronized
specimen traversal; the concrete interface; forward and converse quantitative
robustness conclusions for certified initial fillings. In particular,
determinism is not used to infer reverse padding: the already proved
match-reflection theorem supplies that essential direction.

There are two boundaries to this milestone. First, correspondence between
these coordinate certificates and the paper's explicitly focused graphs,
unique-hole inventories and replacement tuples remains to be proved.
The addressed-closing-context connection was open at this milestone;
Section~\ref{sec:copying-contextual-robustness} subsequently constructs those
certificates for every original closing context and extends initial lifting
to that entire observation class.
Second, the concrete template families still require their output-closure,
input-safety and forward/converse robustness proofs. The generic results below
retain those application conditions and do not assert Proposition~12.2.
Appendix~B.1 removes the completed interface and locality obligations while
retaining these correspondence and template tasks. A cautious estimate is
60--70\% of the total mathematical effort, with low confidence.

\Needspace{12\baselineskip}
\subsection{Retaining the actual ordered interactions}

\termdef{copying-retained-interaction}
A \emph{retained copying interaction} for an edge $e$ of a context $C$ is a
certificate that its replacement component has just one edge, whose source
list is the complete ordered input boundary and whose target list is the
complete ordered output boundary. Its label is a specified primitive exactly
when the original label is that primitive. This permits rebuilding a box's
body while keeping the box opaque to outer search. It does not assume a
search simulation. The position map sends a token $(m,v)$ to $(m,h(v))$, where
$h$ is the expansion's canonical host embedding.

\lean{Copying/ExpansionSearch}{12}{17}{Search.Expansion.Retained}

\lean{Copying/ExpansionSearch}{19}{19}{Search.Expansion.position}

For a retained interaction, the expanded edge's source and target lists
are precisely the original lists mapped by the host embedding. Substituting
the certificate's complete boundary lists into the expansion incidence
formula proves both equalities. These are ordered-list equalities, so they
preserve the argument index used by returning search.

\lean{Copying/ExpansionSearch}{21}{26}{Search.Expansion.Retained.source_eq}

\lean{Copying/ExpansionSearch}{28}{33}{Search.Expansion.Retained.target_eq}

The preservation condition requires a retained certificate for each
interactive primitive edge of the context. The origin condition requires
every interactive edge of a replacement component either to come from such
a retained copy or to replace an original hole. These two structural
conditions have different roles: preservation lifts context search, whereas
origins account for every possible filled-graph search interaction.

\lean{Copying/ExpansionSearch}{36}{37}{Search.Expansion.Preserves}

\lean{Copying/ExpansionSearch}{40}{42}{Search.Expansion.Origins}

Every context search step lifts under preservation. In an entering step,
the retained source equation supplies the same source attachment and the
entering function commutes with its vertex map. In a backtracking step, the
target equation retains the same ordered argument index, while the source
equation retains the operation head. Naturality of the resume function then
gives exactly the mapped successor.

\lean{Copying/ExpansionSearch}{44}{58}{Search.Expansion.step}

A position is entering when its vertex belongs to a hole's source list,
and exiting when it belongs to a hole's target list. The relevant boundary
predicate is mode-sensitive: entering in search mode, or exiting in
backtracking mode. These predicates inspect the actual original hole edge;
they do not require distinct hole names or injective attachments.

\lean{Copying/ExpansionSearch}{60}{60}{Search.Expansion.Entering}

\lean{Copying/ExpansionSearch}{61}{61}{Search.Expansion.Exiting}

\lean{Copying/ExpansionSearch}{62}{62}{Search.Expansion.Boundary}

If an expanded source or target incidence meets an old host vertex, the
owning component's original edge has the corresponding incidence at that
vertex. The proof splits the local vertex into input, output and internal
cases. Fresh internal tags cannot equal an old host tag; the two boundary
cases reduce to the original port maps. These ownership lemmas also apply
when a hole's source and target are attached to the same host vertex.

\lean{Copying/ExpansionSearch}{64}{70}{Search.Expansion.source_owner}

\lean{Copying/ExpansionSearch}{72}{78}{Search.Expansion.target_owner}

An actual search step starting at a host position either reflects to an
original context step, with exactly the mapped successor, or starts at the
mode-sensitive hole boundary. To prove this, classify the interacting edge
by its origin. A hole origin gives the boundary alternative by ownership.
A retained origin has the original primitive label and ordered incidences.
Injectivity of the host map recovers the source membership, or the singleton
source and the returning argument index, and hence the original rule.
Combining this reflection with preservation gives an exact one-step
equivalence away from the boundary.

\lean{Copying/ExpansionSearch}{82}{127}{Search.Expansion.reflect_or_boundary}

\lean{Copying/ExpansionSearch}{129}{134}{Search.Expansion.step_iff}

\Needspace{12\baselineskip}
\subsection{One-wayness and the jointly rooted boundary obstruction}

\termdef{copying-one-way}
A copying hypernet is \emph{one-way} when none of its term outputs is
operation-reachable from an input. The reachability predicate was obtained
from the actual search-history invariant. The path correspondence below
shows that this is exactly the paper's absence of an operation-only path
from a term input to a term output, rather than an additional restriction
on atom or thunk inputs.

\lean{Copying/OneWay}{8}{10}{OneWay}

A component vertex which is not an input cannot become a global input
after insertion. Target coverage gives it an incoming component edge, and
that incidence survives insertion; global inputs have no incoming target.
This fact supplies the base obstruction in the path-extraction proof.

\lean{Copying/OneWay}{18}{21}{Search.Expansion.insert_not_input}

Suppose an inserted vertex is operation-reachable in the expanded graph.
Then every component vertex representing that occurrence is operation-reachable
inside its component. Induct on global reachability. A local input already
has a trivial path. Otherwise target coverage gives a local predecessor;
target uniqueness identifies it with the final global operation edge.
The operation's singleton source identifies the preceding occurrence, to
which the induction hypothesis applies. This proof needs neither global
acyclicity nor an injective hole attachment.

Consequently, attaching a term output of a one-way component to a host
vertex prevents operation reachability at that host occurrence: extraction
would contradict the component's one-way property.

\lean{Copying/OneWay}{25}{55}{Search.Expansion.extract_at}

\lean{Copying/OneWay}{58}{65}{Search.Expansion.oneWay_target}

For paired replacements, output closure requires that at every original
hole at least one of the two components is one-way. The choice may differ
from hole to hole. If both filled tokens are rooted and their common host
vertex has term type, that vertex cannot be a hole target. Indeed, rootedness
gives operation reachability on both sides; use whichever side is one-way
at the selected hole to obtain a contradiction.

\lean{Copying/OneWay}{68}{70}{Search.Expansion.OutputClosedPair}

\lean{Copying/OneWay}{72}{79}{Search.Expansion.not_exiting}

A settled token is either backtracking or searching away from hole
inputs. For jointly rooted paired fillings, it is away from the entire
mode-sensitive boundary: settledness excludes search entry, and the preceding
obstruction excludes backtracking exit. Independently, a context search step
always ends remote. If it ends in rewrite mode, its vertex is the source
of a genuine operation or term contraction; source uniqueness prevents it
from being the source of a hole as well. All other modes are remote by
definition.

\lean{Copying/OneWay}{81}{90}{Search.Expansion.settled_away}

\lean{Copying/OneWay}{93}{122}{Search.Expansion.remote_target}

\Needspace{10\baselineskip}
An operation label in the recursive alphabet is exactly a primitive
operation label. Case analysis on depth and on the label proves the
equivalence. Induction then identifies the reachability predicate with an
actual reverse operation trace to an input: a reachability constructor adds
the corresponding operation edge, and a trace constructor recovers the
operation's singleton source and the preceding reachability proof.

\lean{Copying/OneWayPaths}{7}{20}{isOperation_iff}

\lean{Copying/OneWayPaths}{24}{41}{Search.operationReachable_iff}

Every nonempty operation trace ending at an external input has a term
input there, since the last operation's unique source has term type. The
proof follows the trace to its final edge. An output-to-input trace cannot
be empty, because these boundary constructors are disjoint. Reversing the
edge list therefore proves the promised equivalence: one-wayness holds
exactly when there is no forward operation path from a term input to a
term output, at any interface or depth.

\lean{Copying/OneWayPaths}{45}{62}{Search.operationTrace_input_term}

\lean{Copying/OneWayPaths}{66}{77}{oneWay_iff_paths}

\Needspace{12\baselineskip}
\subsection{Related replacements inside original boxes}

A copying pre-template is a family of relations on actual hypernets,
indexed by depth and the two ordered interfaces. Related graphs must both
be globally focus-free, including all nested bodies.
\termdef{copying-output-closed-template}
An \emph{output-closed copying pre-template} requires each related pair to
have a one-way member. This condition is on proposed template instances;
it is not an assumed law of the copying machine.

\lean{Copying/DeepFillings}{9}{11}{PreTemplate}

\lean{Copying/DeepFillings}{13}{14}{PreTemplate.OutputClosed}

A label copy has one specified edge with a literal specified label,
complete ordered boundary incidences, and no other edges. The canonical
corolla supplies such a certificate directly. If the label is recursively
focus-free, the whole graph is focus-free, because every edge is the
specified edge. A label copy of an original context edge also supplies a
retained-interaction certificate by its label equality.

\lean{Copying/DeepFillings}{17}{22}{LabelCopy}

\lean{Copying/DeepFillings}{24}{26}{LabelCopy.corolla}

\lean{Copying/DeepFillings}{28}{30}{LabelCopy.focusFree}

\lean{Copying/DeepFillings}{32}{38}{LabelCopy.retained}

A label copy of a non-operation is one-way. Every term output is a
target of its sole edge. Operation reachability at such a target would
contradict the earlier non-operation-target obstruction. This argument is
useful even for retained holes, which are non-operation labels.

\lean{Copying/DeepFillings}{40}{45}{LabelCopy.oneWay}

\termdef{copying-recursive-filling}
A \emph{recursive paired copying filling} expands a common context by two
families of edge components, with a certificate at each original edge.
The certificate has three constructors. A retained edge is replaced on both
sides by a literal label copy. An original hole is replaced by an actual
pre-template pair. An original box is replaced by two complete box edges
with the same bound-variable, free-variable and free-atom interfaces, whose
bodies are recursively paired fillings of the original body.

The two mutually defined predicates express these graph equations and edge
certificates. Recursion descends through original boxes only: a newly inserted
hole replacement is not traversed again. At this stage distinct names and
exact original-hole tuples are not part of the certificate. Those paper
correspondence obligations remain explicit.

\lean{Copying/DeepFillings}{48}{53}{DeepFilling}

\lean{Copying/DeepFillings}{56}{70}{DeepFillingEdge}

If the original context is globally focus-free, both recursively filled
graphs are focus-free. Induct on depth. Retained copies use the original
label property; actual hole replacements use the pre-template property;
rebuilt boxes use the induction hypothesis for their complete bodies.
Expansion then assembles these local properties at every edge.

\lean{Copying/DeepFillings}{74}{96}{DeepFilling.focusFree}

A rebuilt box retains outer search interactions even when its body
changes: neither old nor new box label is a primitive label, and the full
boundary lists remain fixed. It follows by cases on an edge certificate
that every component pair either has retained-interaction certificates on
both sides or is a related replacement at an actual original hole.

\lean{Copying/DeepFillings}{99}{106}{LabelCopy.retainedBox}

\lean{Copying/DeepFillings}{109}{116}{DeepFillingEdge.origin}

The structural certificates now imply both search-preservation conditions
and both origin conditions. An interactive original primitive cannot be a
hole; all its certified copies must therefore be retained. Conversely, the
origin classification accounts for every interactive inserted edge.
Output closure of the pre-template gives paired output closure at all
original holes. A genuinely filled hole uses its related pair; a retained
hole is itself one-way by the non-operation argument; a box cannot be a
hole. No operational simulation is included as certificate data.

\lean{Copying/DeepFillings}{119}{134}{deepFilling_search}

\lean{Copying/DeepFillings}{138}{144}{deepFilling_outputClosed}

\Needspace{12\baselineskip}
\subsection{An exact machine on representation classes}

The representation relation is the equivalence relation generated by
recursive state isomorphism and one level of depth padding. It includes
reflexivity, symmetry and transitivity explicitly. These generators are
structural representation changes, not additional machine transitions.
Finality and search rootedness are invariant under every generator and
hence under their equivalence closure, by induction on its derivation.

\lean{Copying/StateClasses}{9}{14}{Combined.Representation}

\lean{Copying/StateClasses}{18}{24}{Combined.Representation.final_iff}

\lean{Copying/StateClasses}{26}{34}{Combined.Representation.rooted_iff}

Every representation equivalence is an exact one-step bisimulation.
For recursive isomorphism, use the existing full-machine transport in both
directions. For padding, forward transport is already proved. In the reverse
direction apply the earlier exact padded-run reflection theorem to the
one-step run; its reflected first step supplies the required old successor,
and its representation equation aligns the results. Reflexivity is immediate;
symmetry exchanges the two directions; transitivity composes witnesses.
The transport lemma exposes the forward half of this result.

This proof depends essentially on reverse matching and run reflection.
Successor determinism alone would not justify eliminating padding.

\lean{Copying/StateClasses}{38}{73}{Combined.Representation.bisimulation}

\lean{Copying/StateClasses}{75}{76}{Combined.Representation.transport}

\termdef{copying-representation-class}
A \emph{copying representation class} is the quotient of unbounded states
by this generated equivalence relation. The setoid packages its three
equivalence laws, and the constructor maps each actual state to its class.
Finality and rootedness descend to the quotient using the invariance proofs;
their truth does not depend on the chosen representative.

\lean{Copying/StateClasses}{79}{81}{Combined.representationSetoid}

\lean{Copying/StateClasses}{83}{83}{Combined.StateClass}

\lean{Copying/StateClasses}{86}{86}{Combined.StateClass.ofState}

\lean{Copying/StateClasses}{88}{89}{Combined.StateClass.Final}

\lean{Copying/StateClasses}{90}{90}{Combined.StateClass.Rooted}

A class step means that some representatives are related by exactly one
original unbounded-machine step. Every original step therefore projects.
Conversely, given a class step and any chosen source representative, the
bisimulation transports its witnessed step back to that representative,
producing an actual successor in the required target class. Thus changing
representatives cannot enable a transition unavailable at the chosen source.

\lean{Copying/StateClasses}{94}{94}{Combined.StateClass.Step}

\lean{Copying/StateClasses}{96}{96}{Combined.StateClass.project}

\lean{Copying/StateClasses}{99}{103}{Combined.StateClass.reflect}

To prove literal determinism of the class relation, reflect two class
steps to the same original source. Existing successor coherence makes their
actual successors recursively isomorphic, so the quotient identifies them.
A final class has no successor, because reflecting one would contradict
original final normality. These facts construct a genuine abstract machine
and discharge its deterministic-machine property without any additional
semantic premise.

\lean{Copying/StateClasses}{105}{109}{Combined.StateClass.deterministic}

\lean{Copying/StateClasses}{111}{114}{Combined.StateClass.final_normal}

\lean{Copying/StateClasses}{116}{119}{Combined.StateClass.machine}

\lean{Copying/StateClasses}{121}{121}{Combined.StateClass.machine_deterministic}

\Needspace{12\baselineskip}
Rootedness is preserved by class steps: choose the witnessed actual
transition and apply full unbounded refocusing. Induction projects each
original run with its exact length. The reverse induction first reflects
one class step from the specified source, then continues from its actual
successor. Consequently every class run has an actual representative run
of the same length from any fixed original starting state, and therefore
at that state's original finite allowance.

\lean{Copying/StateClasses}{123}{126}{Combined.StateClass.rooted_step}

\lean{Copying/StateClasses}{128}{132}{Combined.StateClass.project_run}

\lean{Copying/StateClasses}{136}{145}{Combined.StateClass.reflect_run}

A class has a successful run of length $n$ if and only if its chosen
original representative does. Project or reflect the run and use finality
invariance. Applying these exact correspondences to both sides of quantitative
refinement proves its equivalence for every cost relation $Q$. The class
execution initializes by projecting the unchanged original program state.
Padding contributes no counted transition and introduces no new successful
cost.

\lean{Copying/StateClasses}{147}{154}{Combined.StateClass.successful_iff}

\lean{Copying/StateClasses}{157}{167}{Combined.StateClass.refines_iff}

\lean{Copying/StateClasses}{170}{172}{Combined.StateClass.execution}

\Needspace{12\baselineskip}
\subsection{The concrete specimen and its locality law}

A full transition from a non-rewrite token is necessarily an actual
search step on the unchanged graph. This follows by cases on the combined
transition: the contextual rewrite constructor requires rewrite mode, so
cannot occur. The statement extracts both the genuine search witness and
the exact successor-state equation.

\lean{Copying/Specimens}{9}{14}{Combined.nonrewrite_search}

\termdef{copying-specimen}
A \emph{copying coordinate specimen} consists of a depth, a focus-free context
with sole term input and no outputs, two edge-component families with recursive
paired-filling certificates, hole-freedom proofs for both expanded graphs,
and a term-typed token on the original context. It also records membership
in an explicit predicate of allowed contexts. Its left and right states
are the actual unbounded filled graphs at the host image of the token.
Its entering predicate is the original context's hole-source predicate.

\lean{Copying/Specimens}{19}{32}{Specimen}

\lean{Copying/Specimens}{37}{37}{Specimen.leftState}

\lean{Copying/Specimens}{38}{38}{Specimen.rightState}

\lean{Copying/Specimens}{39}{39}{Specimen.Entering}

\Needspace{12\baselineskip}
Both filled graphs are globally focus-free by the recursive-filling
theorem. Relocating a specimen changes only its token and typing proof;
the context, its allowed-membership proof, every component and every
filling certificate remain the same. The two filled states agree on finality:
their modes coincide, and the host embedding preserves exactly which
vertices are external inputs.

\lean{Copying/Specimens}{41}{42}{Specimen.filled_free}

\lean{Copying/Specimens}{45}{46}{Specimen.relocate}

\lean{Copying/Specimens}{48}{51}{Specimen.same_final}

\termdef{copying-specimen-locality}
\emph{Copying specimen locality} says that if an output-closed specimen is
jointly rooted and settled, and either actual state can step, then both
states can take one full-machine step to the two sides of a remote successor
specimen. Derive preservation, origins and paired output closure from the
filling certificates. Settledness and joint rootedness place the token away
from the hole boundary. Its mode is not rewrite, so the enabled transition
is search. Boundary reflection recovers one actual context step; preservation
lifts that same step to both filled graphs. Its target has term type and is
remote, so relocating the original specimen constructs the required successor.

This proves synchronized search in the full machine; rewriting at an entering
request is handled by template safety and robustness, not by this locality
lemma. Since relocation keeps the context unchanged, no plugging-closure
assumption is needed for this step. That observation does not settle the
separate global context-closure audit in Appendix~B.

\lean{Copying/Specimens}{55}{80}{Specimen.context_step}

The initial specimen relocates the token to search at the sole input.
Both filled sides are rooted by the zero-step search witness. A class step
from a specified non-rewrite original state can also be reflected to actual
search at that exact depth: first reflect the class step, then use the
non-rewrite extraction lemma. This additional statement makes the fixed-source
interpretation explicit.

\lean{Copying/Specimens}{83}{83}{Specimen.initial}

\lean{Copying/Specimens}{85}{86}{Specimen.initial_rooted}

\lean{Copying/Specimens}{91}{99}{Combined.StateClass.search_of_nonrewrite}

The concrete specimen system uses representation classes of the actual
left and right states, the original token mode and entering predicate, and
the descended rootedness predicate. Full refocusing proves its rooted-step
field; common host-boundary classification proves its finality field.
For its context-step field, reflect the enabled class step to the chosen
original state, apply concrete locality, and project both resulting steps.
Thus every abstract interface field is proved. Output closure remains the
pre-template application condition. Any certified initial specimen supplies
the abstract lifting witness because it is remote and jointly rooted.

\lean{Copying/Specimens}{107}{123}{Specimens.system}

\lean{Copying/Specimens}{126}{129}{Specimens.initial_lifting}

Apply the abstract forward sufficiency theorem to this concrete system
and its proved deterministic machine. The quantitative class-refinement
correspondence transfers its conclusion back to actual unbounded states.
The only remaining premises are the proposed template's input safety,
robustness, the reasonable cost triple, and the specified lifting relation.
The converse theorem uses robustness of the opposite system and the original
converse-cost safety condition. No determinism, refocusing or context-locality
law is left as an application premise.

\lean{Copying/Specimens}{133}{139}{Specimens.sufficiency}

\lean{Copying/Specimens}{142}{148}{Specimens.sufficiency_converse}

\Needspace{12\baselineskip}
\subsection{The original observations and the exact initial-filling conclusion}

\termdef{copying-observation-correspondence}
The \emph{copying class observation correspondence} identifies the old and
new executions on all existing observations. First, successful-run equivalence
gives termination equivalence. Then quantify state refinement over the same
original binding-free closing contexts to obtain contextual refinement and
equivalence. Finally quantify over the original terms' adequate depth
allowances and use their unchanged translations. This proves equivalence
of the two term-observation predicates without weakening their context class
or assuming termination.

\lean{Copying/ClassObservations}{7}{8}{Combined.StateClass.terminates_iff}

\lean{Copying/ClassObservations}{10}{12}{Combined.StateClass.contextual_refines_iff}

\lean{Copying/ClassObservations}{14}{16}{Combined.StateClass.contextual_equiv_iff}

\lean{Copying/ClassObservations}{20}{23}{Combined.StateClass.observationalEquiv_iff}

A state and its padded state have equal classes, directly by the padding
generator. More generally, any representation equivalence preserves successful
execution at each exact cost, by quotient equality and successful-run
correspondence. This includes arbitrary finite alternations of padding,
its reverse, and recursive renaming.

\lean{Copying/ClassObservations}{26}{27}{Combined.StateClass.padding_eq}

\lean{Copying/ClassObservations}{31}{35}{Combined.StateClass.representation_successful_iff}

Given an allowed focus-free context and a recursive paired filling with
hole-free results $G,H$, construct its initial coordinate specimen. Output
closure, both input-safety clauses, both robustness clauses and a reasonable
cost triple then imply quantitative equivalence of the actual initial
unbounded states of $G$ and $H$. The proof composes initial lifting with the
two concrete sufficiency theorems.

This is the exact scope of the current initial-filling result. The observation
correspondences above already range over all original closing contexts, but
using this result for each of them still requires constructing their paired
filling certificates. It is not a theorem that every proposed graph relation
is observationally equivalent, and no one of the eleven required template
families has been discharged by the generic implication.

\lean{Copying/ClassObservations}{43}{57}{Specimens.equivalence_of_filling}

\Needspace{12\baselineskip}
\subsection{Concrete boundary, box and cost regressions}

The regression pre-template relates two numeric corollas at a common
depth. Its relation records the actual natural numbers and graph equalities;
both sides are focus-free. It is output-closed because their output interfaces
are empty. This relation tests locality only: no safety or robustness claim
is made for arbitrary pairs of different numbers.

\lean{Copying/SpecimenExamples}{10}{11}{SpecimenExamples.NumberPair}

\lean{Copying/SpecimenExamples}{13}{15}{SpecimenExamples.numbers}

\lean{Copying/SpecimenExamples}{17}{19}{SpecimenExamples.numbers_closed}

\Needspace{16\baselineskip}
The outer example context is unary negation applied to an actual term
hole. Its replacement family retains the negation edge and inserts the
chosen numeric corolla at the hole. These choices give the recursive
edge certificates for a pair of numbers and construct a genuine specimen,
with focus freedom, hole freedom and token typing checked on the actual
expanded graphs.

\lean{Copying/SpecimenExamples}{21}{21}{SpecimenExamples.hole}

\lean{Copying/SpecimenExamples}{22}{23}{SpecimenExamples.context}

\lean{Copying/SpecimenExamples}{25}{31}{SpecimenExamples.family}

\lean{Copying/SpecimenExamples}{33}{37}{SpecimenExamples.paired}

\lean{Copying/SpecimenExamples}{39}{50}{SpecimenExamples.specimen}

The internal argument is the serial attachment between the operation
and the hole. Direct incidence calculations prove the context's entering
step to this argument and its returning step from backtracking there to
rewrite at the outer head. The entry target is an actual hole boundary;
this confirms why unrestricted search reflection would be incorrect.

\lean{Copying/SpecimenExamples}{52}{53}{SpecimenExamples.argument}

\lean{Copying/SpecimenExamples}{55}{58}{SpecimenExamples.context_enter}

\lean{Copying/SpecimenExamples}{60}{61}{SpecimenExamples.context_resume}

\lean{Copying/SpecimenExamples}{64}{65}{SpecimenExamples.entry_crosses_boundary}

Both numeric fillings have a two-search-step path to backtracking at
the internal argument: enter the outer operation, then return from its
numeric value. Thus the returned pair is jointly rooted. Applying the
proved specimen-locality theorem yields synchronized full-machine steps
to a remote specimen. The local return is not globally final, because the
internal attachment is not the external input.

\lean{Copying/SpecimenExamples}{68}{79}{SpecimenExamples.returned_rooted}

\lean{Copying/SpecimenExamples}{83}{92}{SpecimenExamples.synchronized_return}

\lean{Copying/SpecimenExamples}{94}{96}{SpecimenExamples.local_return_not_final}

Induction on the padding count identifies every padded initial program
with its original representation class. The existing unchanged translated
application execution then gives a successful class run of exactly seven
steps at every finite padding, by exact successful-run correspondence.
This checks costs on a real application, rather than only on structural
corollas.

\lean{Copying/SpecimenExamples}{99}{107}{SpecimenExamples.initial_padding_class}

\lean{Copying/SpecimenExamples}{111}{115}{SpecimenExamples.application_class_cost}

Unary negation's operation graph is not one-way: its source input reaches
its term output through the operation itself. A term weakening graph, which
also has a term output, is one-way because its incoming edge is not an
operation. These positive and negative cases show that the predicate is
not merely using the absence of outputs in the numeric regression.

\lean{Copying/SpecimenExamples}{117}{122}{SpecimenExamples.negation_not_oneWay}

\lean{Copying/SpecimenExamples}{125}{126}{SpecimenExamples.weakening_oneWay}

\Needspace{23\baselineskip}
At depth zero, fill a literal term-hole body with each numeric corolla;
the family and its recursive filling witness are explicit. Enclose the
original body in a box with zero bound and free variables and atoms. The
outer families are literal copies of the resulting complete box labels,
and the boxed constructor supplies their recursive filling certificate.
The two bodies may differ. Outer search has no step through either filled
box, since a box label is not a primitive search label. This checks the
body-rebuilding case and its opacity separately from the translated
application-cost regression.

\lean{Copying/SpecimenExamples}{128}{128}{SpecimenExamples.bodyContext}

\lean{Copying/SpecimenExamples}{129}{129}{SpecimenExamples.bodyFamily}

\lean{Copying/SpecimenExamples}{131}{133}{SpecimenExamples.body_filling}

\lean{Copying/SpecimenExamples}{135}{135}{SpecimenExamples.boxedContext}

\lean{Copying/SpecimenExamples}{136}{137}{SpecimenExamples.boxedFamily}

\lean{Copying/SpecimenExamples}{141}{145}{SpecimenExamples.boxed_filling}

\lean{Copying/SpecimenExamples}{147}{152}{SpecimenExamples.boxed_search_opaque}

The complete build passes 460 jobs. All 2458 axiom audits pass: 6 abstract,
477 graph and 1975 copying declarations, including all 112 declarations in
this chapter. Dependencies are confined to \texttt{propext},
\texttt{Classical.choice} and \texttt{Quot.sound}; no project-specific axiom
or omitted proof is introduced. The focused/tuple/context correspondence,
eleven-template argument, semantic weakening and final Figure~22 chain remain
the main work toward Proposition~12.2.
